% ============================
% LaTeX Template - Problem Set
% ============================
\documentclass[11pt]{article}

% ---------- Page + layout ----------
\usepackage[margin=1in]{geometry}
\usepackage{setspace}
\setstretch{1.1}
\usepackage{parskip}

% ---------- Math ----------
\usepackage{amsmath, amssymb, amsfonts, amsthm, mathtools}

% ---------- Tables ----------
\usepackage{booktabs}

% ---------- Links ----------
\usepackage[hidelinks]{hyperref}

% ---------- Header / footer ----------
\usepackage{fancyhdr}
\pagestyle{fancy}
\fancyhf{}
\lhead{\textbf{\course}}
\chead{\textbf{\pset}}
\rhead{\textbf{\name}}
\cfoot{\thepage}

% =========================
% Metadata (edit if needed)
% =========================
\newcommand{\name}{Tate Mason}
\newcommand{\course}{Econometrics II}
\newcommand{\pset}{Problem Set 6}
\newcommand{\duedate}{March 25, 2026}

% =========================
% Question / solution macros
% =========================
\newcounter{q}
\renewcommand{\theq}{\arabic{q}}

\newcommand{\question}[1]{%
  \refstepcounter{q}
  \vspace{0.75em}
  \noindent{\Large\textbf{Question \theq.} \normalsize\textbf{#1}}%
  \vspace{0.25em}
  \par
}

\newcounter{subq}[q]
\renewcommand{\thesubq}{\alph{subq}}

\newcommand{\subquestion}[1]{%
  \refstepcounter{subq}
  \vspace{0.35em}
  \noindent\textbf{(\thesubq)} \textbf{#1}\par
}

\newenvironment{solution}{%
  \vspace{0.15em}
  \noindent\textit{Solution. }\ignorespaces
}{\vspace{0.5em}\par}

% =========================
% Easy math macros
% =========================
\newcommand{\R}{\mathbb{R}}
\newcommand{\N}{\mathbb{N}}
\DeclareMathOperator{\E}{\mathbb{E}}
\DeclareMathOperator{\Var}{Var}
\DeclareMathOperator{\Cov}{Cov}
\DeclareMathOperator{\plim}{plim}
\DeclareMathOperator*{\argmin}{argmin}
\DeclareMathOperator*{\argmax}{argmax}
\newcommand{\P}{\mathbb{P}}

\newcommand{\vct}[1]{\mathbf{#1}}
\newcommand{\mtx}[1]{\mathbf{#1}}
\newcommand{\dd}{\,\mathrm{d}}
\newcommand{\eps}{\varepsilon}
\newcommand{\1}{\mathbf{1}}

% =========================
% Document
% =========================
\begin{document}

% ---------- Title block ----------
{\Large\textbf{\course} \hfill \textbf{\pset}}\\
\textbf{Name:} \name \hfill \textbf{Due:} \duedate\\
\rule{\textwidth}{0.4pt}
\question{Use ``dataHW6 Problem1.mat'' to complete this exercise. The purpose of this problem is to estimate a simple dynamic discrete choice problem. The data contain information on the choices (the $Y$'s), characteristics that do not vary over time (the $X_1$'s), and characteristics that do vary over time (the $X_{1t}$'s). The context of the problem is a dynamic model of labor supply. $Y$ takes values $\{0, 1, 2, 3\}$ reflecting retirement, flex-time, part-time, or full-time participation. The first column $X_1$ allows for choice specific constants, while the second column of $X_1$ is an indicator for gender. $X_{1t}$ is a time-varying measure of health.

    Both sets of $X$'s are individual-level characteristics so the coefficients on these variables need to vary by choice. There is a cost to switching labor supply in the future that is \textit{common across all switches}. The choice in period 0 is given by $LY$. Once an individual retires (or chooses option 0) then that individual makes no further choices (you can verify this as if $Y_{it} = 0$ then $Y_{it+1} = 0$). The columns of $Y$ and $X_{1t}$ indicate observations over time, ten periods in all (so in period 10 it is a static problem). The only uncertainty from the individual's perspective is on their error term. In other words, individuals have perfect foresight about their future health. The error structure is Type 1 extreme value.}

\subquestion{Estimate the dynamic discrete choice model including the discount factor.}

\begin{solution}
  \begin{table}[htbp]
  \centering
    \caption{DDC Estimates}
    \begin{tabular}{lc}
      \hline
      & Parameter \\
      \hline
      \multicolumn{3}{l}{\textit{Flex-Time}} \\
      $\beta_{0,1}$    & -0.961 \\
      $\beta_{Gender,1}$   & 0.368 \\
      $\beta_{Health,1}$ & 0.517 \\

      \multicolumn{3}{l}{\textit{Part-Time}} \\
      $\beta_{0,2}$    & --.051 \\
      $\beta_{Gender,2}$   & 0.687 \\
      $\beta_{Health, 2}$ & 0.759 \\

      \multicolumn{3}{l}{Full-Time} \\
      $\beta_{0,3}$ & -0.412 \\
      $\beta_{Gender,3}$ & 1.029 \\
      $\beta_{Health,3}$ & 0.992 \\
      \hline

      $\delta$ & 0.376 \\
      $\beta$ & 0.896 \\
      \hline
    \end{tabular}
  \end{table}
\end{solution}

\clearpage

\subquestion{What variation in the data allows the discount factor to be identified?}

The discount factor is identified via the health state. It is the only time variant characteristic in $X_{1t}$, so it will pin down $\beta$.

\subquestion{What variation in the data allows the switching cost to be identified?}

Switching cost is identified via the constant, the only characteristic which is choice specific in $X1$. Gender is not something which we are assuming can change in this model, thus it must be the constant identifying the cost $\delta$.
\end{document}
